Target Wake Time (TWT), introduced in IEEE 802.11ax, lets an access point
schedule exactly when each station wakes, transmits, and dozes. Existing TWT
schedulers optimize energy or throughput, treating information freshness at
best as a constraint and offering no performance guarantees. We design the
wake schedule itself for freshness: minimize the weighted average Age of
Information (AoI) over stations subject to per-station energy budgets, where
the decision variables are the TWT triples (wake interval, offset, service
period duration). We derive a renewal-exact AoI model for TWT under per-SP
block fading and validate it against packet-level 802.11ax simulation with
${\sim}1\%$ mean error. We show that, unlike preemptive scheduling, non-preemptive TWT
packing can be infeasible at schedule density~1, and identify the granularity
condition under which a small-first best-fit packer provably succeeds. Around
this we build \textsc{Harmonic-Greedy}, a scheduler combining a convex
relaxation, anchor-optimized power-of-two rounding, and a best-of-uniform
safeguard, and prove it is a constant-factor approximation: $4/\ln 2 \approx
5.77$ under a mild granularity assumption and $6/\ln 2 \approx 8.66$
unconditionally. We implement the complete system in ns-3 -- a TWT wake/doze
mechanism integrated with the power-save architecture, plus the scheduler --
and show that it is the only scheduler that stays near a relaxation lower
bound across all regimes: against a strong energy-greedy baseline it ties
when per-station energy floors already pin the periods, and wins by
$5$--$36\%$ exactly where the schedule density is binding and must be
redistributed by AoI weight or channel quality rather than by energy budget
-- the regime our analysis identifies.
